% !TEX root = ../main.tex
%% Corrigé du devoir surveillé n° 6 — valeurs vérifiées : verifications/devoirs.py

\corrigeDevoir{6}

\begin{corrigeBac}{7}{La célérité des ultrasons}
\begin{questions}
\item Oui : ce sont des ondes \res{\text{mécaniques}}, qui se propagent dans l'air comme le son. Leur fréquence, \qty{40,0}{kHz}, dépasse
\qty{20}{kHz}, limite supérieure de l'audition humaine. \pts{0,5 pt}
\item L'erreur de repérage de la position du récepteur est commise une seule fois pour dix longueurs d'onde : l'incertitude
sur $\lambda$ est \res{\text{divisée par dix}}. \pts{1 pt}
\item $\overline{d}=\res{\qty{8,519}{cm}}$ ; $s=\res{\qty{0,054}{cm}}$ (calculatrice, $\sigma_{n-1}$). $\lambda=\dfrac{\overline{d}}{10}=\res{\qty{0,8519}{cm}}$. \pts{1,5 pt}
\item $v=\lambda\,f=\num{8,519e-3}\times\num{40,0e3}=\res{\qty{340,8}{m.s^{-1}}}$. \pts{1 pt}
\item $u(\overline{d})=\dfrac{s}{\sqrt{8}}=\qty{0,019}{cm}$ ; $v$ est proportionnelle à $d$ : $u(v)=\dfrac{u(\overline{d})}{10}\times f=\num{1,9e-5}\times\num{40,0e3}\approx\qty{0,8}{m.s^{-1}}$.
Résultat : $v=\res{\qty{340,8}{m.s^{-1}}}$ avec $u(v)=\res{\qty{0,8}{m.s^{-1}}}$. \pts{1,5 pt}
\item À \qty{15}{\degreeCelsius} : $v\approx331+0{,}6\times15=\qty{340}{m.s^{-1}}$ ; l'écart, \qty{0,8}{m.s^{-1}}, est de l'ordre de $u$ : résultat \res{\text{compatible}}.
L'écart à \qty{343}{m.s^{-1}} vaut environ $3\,u$ : cette valeur ne convient pas, car elle correspond à une \res{\text{autre température}}. \pts{1,5 pt}
\end{questions}
\retiens{Avant de comparer à une valeur de référence, vérifier qu'elle correspond bien aux conditions de l'expérience.}
\end{corrigeBac}

\begin{corrigeBac}{7}{L'objectif d'un appareil photo}
\begin{questions}
\item $C=\dfrac{1}{f'}=\dfrac{1}{0{,}050}=\res{\num{20}\,\delta}$. \pts{0,5 pt}
\item $\dfrac{1}{\alg{OA'}}=\dfrac{1}{f'}+\dfrac{1}{\alg{OA}}=\dfrac{1}{0{,}050}+\dfrac{1}{-5{,}00}=\qty{19,8}{m^{-1}}$, d'où $\alg{OA'}=\res{\qty{50,5}{mm}}$. \pts{1,5 pt}
\item $\gamma=\dfrac{\alg{OA'}}{\alg{OA}}=\dfrac{0{,}0505}{-5{,}00}=\res{\num{-0,0101}}$ ; $\alg{A'B'}=\gamma\times\alg{AB}=\res{\qty{-20,2}{mm}}$ : image renversée de
\qty{20,2}{mm}, inférieure à \qty{24}{mm} : la personne \res{\text{tient entière}} sur le capteur. \pts{1,5 pt}
\item Pour le paysage, $\alg{OA'}=f'=\qty{50,0}{mm}$ : il faut \res{\text{rapprocher l'objectif du capteur}} de $50{,}5-50{,}0=\res{\qty{0,5}{mm}}$. \pts{1 pt}
\item Construction ci-dessous. On mesure $\alg{OA'}\approx\qty{7,5}{cm}$ et $\alg{A'B'}\approx\qty{-1,0}{cm}$ (en tenant compte de l'échelle). Vérification :
$\dfrac{1}{\alg{OA'}}=\dfrac15-\dfrac{1}{15}$, d'où $\alg{OA'}=\res{\qty{7,5}{cm}}$ ; $\gamma=\dfrac{7{,}5}{-15}=-0{,}5$ et $\alg{A'B'}=\res{\qty{-1,0}{cm}}$. \pts{2 pts}
\item Le \res{\text{cristallin}} joue le rôle de l'objectif, la \res{\text{rétine}} celui du capteur. \pts{0,5 pt}
\end{questions}
\schema[corrige]{ds6-lentille}
\piege{Oublier le signe : un objet réel a $\alg{OA}<0$ ; ici $\alg{OA}=\qty{-5,00}{m}$, et non \qty{5,00}{m}.}
\end{corrigeBac}

\begin{corrigeBac}{6}{L'écran d'un smartphone}
\begin{questions}
\item Synthèse \res{\text{additive}} : on superpose des lumières colorées émises par les sous-pixels. \pts{0,5 pt}
\item Rouge + vert : \res{\text{jaune}} ; les trois : \res{\text{blanc}} ; vert + bleu : \res{\text{cyan}}. \pts{1 pt}
\item $E=\dfrac{h\,c}{\lambda}$ : rouge $\dfrac{\num{6,63e-34}\times\num{3,00e8}}{\num{625e-9}\times\num{1,60e-19}}=\res{\qty{1,99}{eV}}$ ; vert \res{\qty{2,37}{eV}} ; bleu \res{\qty{2,67}{eV}}. \pts{1,5 pt}
\item $\nu=\dfrac{c}{\lambda}=\dfrac{\num{3,00e8}}{\num{525e-9}}=\res{\qty{5,71e14}{Hz}}$ ; $E=h\,\nu=\res{\qty{3,79e-19}{J}}$. \pts{1 pt}
\item L'écart entre niveaux est égal à l'énergie du photon émis : le plus grand correspond au photon le plus énergétique, celui
de plus courte longueur d'onde : le sous-pixel \res{\text{bleu}} (\qty{2,67}{eV}). \pts{1 pt}
\item Il reste surtout du rouge et du vert, dont la superposition donne du jaune : la page prend une teinte \res{\text{jaune
orangé}}. \pts{1 pt}
\end{questions}
\end{corrigeBac}
